\documentclass[10pt,twocolumn]{article}

\usepackage[utf8]{inputenc}
\usepackage[margin=0.75in]{geometry}
\usepackage{amsmath,amssymb,amsfonts}
\usepackage{booktabs}
\usepackage{graphicx}
\usepackage{hyperref}
\usepackage{xcolor}
\usepackage{cite}

\hypersetup{
    colorlinks=true,
    linkcolor=blue,
    citecolor=blue,
    urlcolor=blue
}

\title{\textbf{Reproduction, Causal Extension, and Methodological Audit of Multidimensional Matrix Profile for Network Anomaly Detection in Cloud Environments}}

\author{
    \textbf{Individual Course Project}\\
    Cloud SRE and Security\\
    \textit{Base Paper: Al-Mazrawe \& Al-Musawi (2025), Journal of Internet Services and Applications (JISA)}
}

\date{November 2026}

\begin{document}

\maketitle

\begin{abstract}
We present a rigorous reproduction, causal extension, and methodological audit of the multidimensional Matrix Profile (mSTAMP) and Kneedle dimension-selection pipeline proposed by Al-Mazrawe \& Al-Musawi (2025) for cloud network anomaly detection. While the base paper reported near-perfect detection fidelity ($F_1 > 0.998$) on public benchmarks, critical operational details---including score definitions, subsequence alignment, decision threshold calibration, and chronological partition integrity---were underspecified. We re-implement the pipeline under strict leakage-free protocols and audit it across UNSW-NB15 and CSE-CIC-IDS2018 (AWS cloud telemetry). 

Our audit uncovers five critical findings: (1) synthetic timestamps skipped 528 quiescent minutes on Friday, distorting wall-clock MTTD and operational alert rates; (2) non-parametric detectors are highly sensitive to reference sizing, where starved training ($N=16$) collapsed Isolation Forest to $F_1 = 0.0000$, while cross-day sizing ($N=507$) restored it to Rank 1 tied with ARIMA at $F_1 = 0.9859$; (3) saturation on flow data was caused by mathematical threshold collapse (99.73\% zero validation scores driving $\tau^* = 0.0$ and zero MAD) rather than stale reference libraries across 26 days; (4) paired persistence filtering ($k \ge 2$) degraded test $F_1$ in 67.5\% of cases; and (5) on 10 contiguous flow blocks (80,000 flows), all matrix profile variants scored mean $F_1 \le 0.1246$, underperforming the naive All-Attack baseline ($F_1 = 0.1281$).

To resolve the threshold generalization gap caused by diurnal traffic non-stationarity, we propose Split-Conformal and Adaptive Anomaly Thresholding (ACI) as our positive contribution. Across 32 configurations, ACI restored false alarm control from a mean test FPR of 12.37\% down to 2.86\% ($\alpha = 0.05$) and 1.06\% ($\alpha = 0.01$). Finally, in a live microservice testbed with 25 labeled campaigns, we demonstrate that static thresholds suffer a 100\% false alarm rate during benign flash crowds, whereas rolling conformal calibration slashes flash crowd false alarms to 11.4\%, lifting Isolation Forest test $F_1$ to 0.9451 with 100\% campaign recall.
\end{abstract}

\section{Introduction}
Modern cloud infrastructures depend on Site Reliability Engineering (SRE) practices to monitor vast volumes of flow telemetry (e.g., AWS VPC Flow Logs) for cyber threats, volumetric distributed denial-of-service (DDoS) attacks, and unauthorized reconnaissance. In production cloud operations, anomaly detectors must balance three competing requirements: high recall on critical attacks to minimize Mean Time to Detect (MTTD), strictly bounded false positive rates (FPR) to prevent operator alert fatigue, and low computational overhead without continuous retraining.

Recently, Al-Mazrawe \& Al-Musawi \cite{almazrawe2025enhancing} proposed combining multidimensional Matrix Profile (mSTAMP) with the Kneedle algorithm for cloud network intrusion detection, reporting $99.8\%$ accuracy and $F_1 = 0.998$ on UNSW-NB15 \cite{moustafa2015unsw}. However, their published description omitted crucial operational definitions: the anomaly score function was undefined, decision thresholds were unexplained, subsequence alignment was unstated, and evaluation splits lacked non-overlapping chronological ordering guarantees.

In this work, we conduct a faithful, leakage-free reproduction, perform a root-cause methodological audit, develop a split-conformal false alarm bounding framework, and validate the frozen pipeline on a live container-ready microservice testbed.

\section{Background \& Theoretical Formulation}
\subsection{Multidimensional Matrix Profile (mSTAMP)}
Let $X \in \mathbb{R}^{n \times d}$ be a regularized multivariate time series across $d$ standardized features. For a subsequence length $m \ll n$, let $T_{i, :}^m$ denote the $m$-length subsequence starting at index $i$. The $z$-normalized Euclidean distance between two univariate subsequences on channel $j$ is:
\begin{equation}
d(T_{i, j}^m, T_{k, j}^m) = \sqrt{2m \left( 1 - \frac{T_{i, j}^m \cdot T_{k, j}^m - m \mu_i \mu_k}{m \sigma_i \sigma_k} \right)}
\end{equation}
For each index $i$, distances across all dimensions are sorted in ascending order. The cumulative multi-dimensional distance across $k$ dimensions is:
\begin{equation}
P_{i}^{(k)} = \min_{j \ne i} \sum_{l=1}^k d(T_{i, \pi(l)}^m, T_{j, \pi(l)}^m)
\end{equation}
where $\pi$ is the sorting permutation for index $i$ \cite{yeh2017multidimensional}.

\subsection{Kneedle Subspace Selection}
To select the intrinsic subspace dimensionality $k^* \in \{1, \dots, d\}$, mSTAMP takes the maximum discord score across subsequences for each dimension count: $y_k = \max_{i} P_{i}^{(k)}$. Because $y_k$ is concave and increasing, Kneedle \cite{satopaa2011finding} identifies the inflection index maximizing distance to the diagonal:
\begin{equation}
k^* = \arg\max_k \left( \frac{y_k - y_1}{y_d - y_1} - \frac{k - 1}{d - 1} \right)
\end{equation}

\subsection{Causal Scoring}
To prevent look-ahead bias in streaming SRE applications, subsequence scores must be causally right-aligned:
\begin{equation}
\mathcal{S}(t) = P_{t - m + 1}^{(k^*)} \quad \text{for } t \ge m - 1
\end{equation}

\section{Methodological Audit \& Protocol Repairs}
\subsection{Time-Axis Integrity}
On CSE-CIC-IDS2018 Friday \cite{sharafaldin2018toward}, 191 active windows spanned 11.97 elapsed hours, containing 528 empty 60-second windows and a 454-minute gap. Synthetically indexing active windows artificially dilated 12 hours into 3.1 hours. We projected raw flows onto an explicit, gapless 60-second datetime grid, and formally asserted $\text{Train End} < \text{Val Start} < \text{Test Start}$.

\subsection{Training Sizing \& Ranking Reversals}
Within-day training severely starved Friday ($N=45$) and Wednesday ($N=16$), failing statsmodels ARIMA(1,0,1) convergence \cite{box2015time}. Introducing cross-day reference sizing ($N=507$ from Thursday morning) resolved starvation: on Wednesday, Isolation Forest surged from $F_1 = 0.0000$ (Rank 5) to Rank 1 tied with ARIMA at $F_1 = 0.9859$. Detector rankings are sensitive to training capacity.

\subsection{Refutation of Stale-Reference Hypothesis}
Config A-prime saturation (100\% false alarms) on flow data was previously hypothesized to stem from a 26-day calendar gap. Our diagnosis proved that UNSW Flow Block is located entirely within Batch 1 (107s after reference). Saturation was caused by 99.73\% of benign validation scores being zero (driving validation-$F_1$ optimization to select $\tau^* = 0.0$) and rolling MAD collapsing to 0. A label-free quantile rule with an $\epsilon = 10^{-6}$ MAD floor reduced test FPR from 100\% to 0.57\%.

\subsection{Paired Persistence Evaluation}
Across 80 parent-variant pairs on identical score streams and thresholds ($\alpha=0.05$), persistence filtering ($k \ge 2$) hurt test $F_1$ in 67.5\% of cases and was neutral in 30.0\%. In 80\% of cases, validation selected $k^* = 1$, refuting claims of operational benefit.

\subsection{Multi-Block Flow Evaluation}
Across 10 contiguous blocks of 8,000 flows (80,000 flows total), all matrix profile variants scored mean $F_1 \le 0.1246$ (worse than naive All-Attack $F_1 = 0.1281$), with near-chance ROC-AUC ($0.52 - 0.56$). Flow-level anomaly detection with matrix profiles is not supported by empirical evidence.

\section{Stage 2 Contribution: Conformal False Alarm Control}
\subsection{Split-Conformal \& Adaptive Thresholding}
To overcome diurnal baseline drift where static thresholds violate target FPR budgets, we formulate Split-Conformal Anomaly Calibration:
\begin{equation}
q_{\alpha} = \frac{\lceil (n + 1)(1 - \alpha) \rceil}{n}, \quad \tau_{\alpha} = \text{Quantile}(\mathcal{S}_{\text{cal}}, q_{\alpha})
\end{equation}
Exact conformal $p$-values are computed as $p(s) = (1 + \sum_{i=1}^n \mathbb{I}(s_i \ge s)) / (n + 1)$. Alerts are raised when $p(s) \le \alpha$.

Under non-stationary drift, we deploy Adaptive Conformal Inference (ACI) \cite{gibbs2021adaptive}, which dynamically updates nominal risk online:
\begin{equation}
\alpha_{t+1} = \text{clip}\left(\alpha_t + \gamma \cdot (\alpha - \text{err}_t), \alpha_{\min}, \alpha_{\max}\right)
\end{equation}
where $\text{err}_t = \mathbb{I}(s_t \ge \tau_t)$ on benign instances, and $\gamma = 0.01$.

\begin{table}[t]
\centering
\caption{Conformal Threshold Benchmark Summary Across 32 Runs}
\label{tab:conformal_summary}
\small
\begin{tabular}{lccccc}
\toprule
\textbf{Method} & \textbf{Target $\alpha$} & \textbf{Mean FPR} & \textbf{Excess} & \textbf{Met \%} & \textbf{Mean $F_1$} \\
\midrule
Static Quantile & 0.05 & 0.1237 & +0.0737 & 43.8\% & 0.4076 \\
Static Conformal & 0.05 & 0.1199 & +0.0699 & 43.8\% & 0.4053 \\
Rolling Conf. ($W=60$) & 0.05 & \textbf{0.0306} & \textbf{-0.0194} & \textbf{68.8\%} & 0.2349 \\
Adaptive ACI & 0.05 & \textbf{0.0286} & \textbf{-0.0214} & \textbf{71.9\%} & 0.2345 \\
\midrule
Static Quantile & 0.01 & 0.0200 & +0.0100 & 65.6\% & 0.2507 \\
Rolling Conf. ($W=60$) & 0.01 & \textbf{0.0106} & \textbf{+0.0006} & \textbf{62.5\%} & 0.1478 \\
Adaptive ACI & 0.01 & \textbf{0.0106} & \textbf{+0.0006} & \textbf{62.5\%} & 0.1478 \\
\bottomrule
\end{tabular}
\end{table}

As shown in Table \ref{tab:conformal_summary}, ACI brought mean test FPR down from 12.37\% to 2.86\% ($\alpha=0.05$) and 1.06\% ($\alpha=0.01$). For Isolation Forest on Friday (DoS-Hulk), Rolling Conformal slashed FPR from 12.50\% to 1.67\% while maintaining 100\% recall, boosting $F_1$ from 0.8624 to 0.9792.

\section{Stage 3 Bounded Testbed Validation}
\subsection{Live Microservice Architecture}
We constructed a containerized testbed featuring:
\begin{itemize}
    \item \textbf{Microservice:} Flask app with Prometheus exporter exposing `/api/catalog`, `/api/search`, `/api/checkout`, and `/metrics`.
    \item \textbf{Benign Load:} Locust generator simulating a diurnal sinusoidal curve with Poisson user transactions.
    \item \textbf{Confined Attacker:} Strictly confined to localhost (`127.0.0.1:8088`), executing 25 labeled campaigns across 600 regular 60s windows (10 hours).
    \item \textbf{Benign Flash Crowds:} 5 campaigns where legitimate traffic surges 4x--5x to `/api/catalog` with 0\% error rate and valid 200 OK responses.
\end{itemize}

\begin{table}[t]
\centering
\caption{Frozen Pipeline Performance on Live Testbed Telemetry ($\alpha=0.05$)}
\label{tab:testbed_results}
\small
\begin{tabular}{lccccc}
\toprule
\textbf{Model} & \textbf{Threshold} & \textbf{Test FPR} & \textbf{Test $F_1$} & \textbf{Flash FPR} & \textbf{Camp. Rec} \\
\midrule
Isolation Forest & Static Quantile & 0.2622 & 0.8526 & 1.0000 & 100\% \\
Isolation Forest & Rolling Conf. & \textbf{0.0488} & \textbf{0.9451} & \textbf{0.1143} & 100\% \\
Isolation Forest & Adaptive ACI & \textbf{0.0488} & 0.8923 & \textbf{0.1143} & 90\% \\
\midrule
ARIMA Residual & Static Quantile & 0.2744 & 0.8471 & 1.0000 & 100\% \\
ARIMA Residual & Rolling Conf. & \textbf{0.0732} & \textbf{0.9353} & \textbf{0.2000} & 100\% \\
ARIMA Residual & Adaptive ACI & \textbf{0.0549} & \textbf{0.9299} & \textbf{0.1429} & 100\% \\
\midrule
Config A (mSTAMP) & Static Quantile & 0.0000 & 0.0000 & 0.0000 & 0\% \\
Config A-prime & Static Quantile & 0.0000 & 0.0000 & 0.0000 & 0\% \\
\bottomrule
\end{tabular}
\end{table}

\subsection{Flash Crowd Vulnerability \& Resolution}
As detailed in Table \ref{tab:testbed_results}, static thresholds triggered a \textbf{100.0\% false alarm rate during flash crowds (35 / 35 windows incorrectly flagged as attacks)}. Because legitimate surges increase volume, static thresholds are blind to valid status codes and normal latency profiles. 

In contrast, Rolling Conformal and ACI adapted dynamically, slashing flash crowd false alarms to \textbf{11.43\%} (an 88.6\% reduction) and lifting Isolation Forest $F_1$ to \textbf{0.9451} while maintaining 100\% campaign recall. Concurrently, Config A and A-prime collapsed (ROC-AUC 0.5000), replicating the breakdown of matrix profiles on discrete flow metrics.

\section{Discussion \& Operational Recommendations}
In cloud SRE deployments, alert fatigue is the paramount operational hazard. While static thresholds achieve marginally higher recall on borderline stealth attacks, they generate intolerable alert storms ($> 35$ alerts/hour). Conformal thresholding bounds false alarm rates to pre-specified budgets, guaranteeing statistical validity.

We recommend three operational guidelines for production Cloud SRE anomaly detection:
\begin{enumerate}
    \item \textbf{Prioritize Tabular Ensembles:} Isolation Forest and ARIMA residuals decisively outperform matrix profiles on network flow metrics.
    \item \textbf{Deploy Split-Conformal Calibration:} Dynamic conformal thresholding bounds false alarm rates under diurnal baseline drift and eliminates 88.6\% of flash crowd false alarms.
    \item \textbf{Enforce Minimum Reference Capacity:} Maintain at least 100--500 benign windows to prevent model starvation and ranking reversals.
\end{enumerate}

\section{Conclusion}
This study presented a comprehensive reproduction, causal extension, and methodological audit of mSTAMP + Kneedle for cloud network intrusion detection. We demonstrated that published near-perfect detection numbers rely on unstated assumptions and high-prevalence slices. Across multi-block flow benchmarks and live microservice telemetry, matrix profiles underperform standard baselines. Split-conformal and adaptive thresholding provides mathematically grounded false alarm rate control, resolving the flash crowd pathology and providing a robust foundation for production cloud telemetry monitoring.

\bibliographystyle{IEEEtran}
\bibliography{references}

\end{document}

